\section{Discussion and conclusions}\label{sec:conclusions}
The metric trial functions of unsymmetric spectral elements suffer from two distinct
deficiencies that grow with the polynomial degree and the element distortion: an
ill-conditioned interpolation \emph{system} and an ill-conditioned interpolation
\emph{problem}. The first is a matter of the basis, the second of the trial space.

\paragraph{Basis.} The trial space, its Lebesgue constant, distortion immunity and the exact
element matrix are basis invariant (\cref{prop:invariance}). The scaled tensor-product Lagrange
basis on GLL points makes the interpolation matrix a permutation on affine elements, so the
unsymmetric element coincides with the conventional one to round-off there (symmetric
$\bm K^e$ for all $p$), and on distorted elements it reduces $\kappa(\bm V)$ to the Lebesgue
function of the interpolation at the reference points (\cref{prop:kappa}). Compared with
monomials it lowers $\kappa$ by $4.6\cdot10^3$--$2.9\cdot10^4$ and the error of the computed shape
functions by $3.6\cdot10^2$--$8.2\cdot10^3$ on the distorted test elements at $p=16$, and it removes the
spurious rejection of elements by a singularity check (five element--degree combinations on the panel meshes). Its
benefit is confined to round-off: approximation, the amplification of nodal perturbations and
the stability of the discrete problem are unchanged, and the static results of the thick cylinder
change only at round-off level. A better basis \emph{cannot} improve the Lebesgue constant.

\paragraph{Hybrid space.} Completeness requires metric functions only for $P_p(\bx)$. The hybrid
space $T_p$ (\cref{def:hybrid}) keeps them for $P_p$ and represents the remaining
$\tfrac12p(p+1)$ functions in parametric coordinates. It retains distortion immunity, frame
invariance and interpolation, reduces to the conventional spectral element on straight-sided
elements, and its interpolant is the conventional one plus the interpolation error of a
polynomial composed with the element map (\cref{prop:structure}), which bounds its Lebesgue
constant by the GLL value plus a term that vanishes for bilinear maps. Unisolvence is equivalent
to the non-singularity of a $d\times d$ matrix $\bm A$ (\cref{prop:factor}). Numerically, the hybrid
space
\begin{itemize}
\item keeps the Lebesgue constant at or near the GLL value on single elements up to $p=16$
  ($\le8.5$, against $5.5\cdot10^4$--$4.0\cdot10^7$ for $Q_p$ in skew coordinates) and $\le95$ on the
  thick-cylinder meshes without folds (against up to $4.2\cdot10^6$);
\item reproduces polynomial fields to $\le10^{-9}$ on all meshes, including folded ones, where
  $Q_p$ in skew coordinates is limited to $10^{-5}$--$10^{-2}$ by round-off;
\item produces no eigenvalues with negative real part and imaginary parts $\le1.8\,\%$ of the
  modulus, where $Q_p$ in skew coordinates produces up to $25$ unstable modes on curved meshes;
\item converges with the optimal order $p$ and quasi-optimally (discrete error $0.78$--$1.00$ times
  the interpolation error) on all non-folded meshes of a mesh sequence with fixed relative
  distortion, on which $Q_p$ in skew coordinates is optimal only for moderate distortion and the
  conventional element converges with $\approx\lfloor p/2\rfloor$; the gain over the conventional
  element reaches $2.6\cdot10^4$ in the energy norm ($p=7$, $\alpha=30^\circ$, $N=8$).
\end{itemize}

\paragraph{Limitations.} (i) No stability theory (discrete inf--sup
condition~\cite{Babuska1971Error}) is available for either unsymmetric element; quasi-optimality
of \PH\ is observed, not proven. (ii) Unisolvence of $T_p$ is guaranteed only near affine maps; it
fails on some folded elements at $p\ge7$, and the complement part of the element matrix has a
singular integrand where $\det\bm J$ changes sign; on folded meshes \PH\ is reliable only for
$p\ge4$--$5$. (iii) The error constant of \PH\ depends on the distortion: where $Q_p$ in skew
coordinates is stable, it can approximate better (by factors up to $5.7$ for the Lam\'e field at
$\alpha=15^\circ$, up to $17$ for $e^x\cos y$ on curved meshes). (iv) The studies are two-dimensional,
linear-elastic and restricted to $p\le16$ (element level) and $p\le8$ (convergence); 3D elements and
non-linear problems remain to be examined.

For high-order unsymmetric spectral elements we therefore recommend the scaled Lagrange basis
for any metric space, and the hybrid space $T_p$ whenever $p\gtrsim4$, the distortion is not
small, or the spectrum matters (transient and modal analysis).

\section*{Reproduction}\label{sec:reproduction}
Code base: USEM-MATLAB after commit \texttt{e293da1} (\texttt{ShapeFunctions2D\_Metric},
\texttt{p.metricBasis}, \texttt{p.metricMethod}). Element and mesh level:
\texttt{studies/metric\_basis/*.m}; patch tests: \texttt{studies/rajendran\_panel};
thick cylinder: \texttt{studies/thick\_cylinder}; unisolvence:
\texttt{studies/manuscript\_metric/matlab/unisolvence\_check.m}; tables and data of this paper:
\texttt{studies/manuscript\_metric/make\_data.py}.
